%
% File:    ./macros/notes.tex
% Author:  Jiří Kučera <sanczes AT gmail.com>
% Date:    2022-09-05 14:18:54 +0200
% Project: Math Notes
% Brief:   Macros for typesetting notes
%
% SPDX-License-Identifier: MIT
%
\input macros/document%
\input macros/drawing%
%
\paperwidth 210 true mm%
\paperheight 297 true mm%
%
\input macros/pdf%
%
% Let @ be a letter:
\makeatletter%
%
% Extra fonts:
%   - title page and cover
\font\lxiisf=NotoSans-Bold-tlf-t1 at 62\Unit%
\font\xxvit=ntx-Italic-tlf-t1 at 25\Unit%
\font\xxiiiit=ntx-Italic-tlf-t1 at 23\Unit%
%
% Text font size switches:
\let\OldNormal\normal%
\def\normal{%
  \OldNormal%
  \abovedisplayskip 10\Unit%
  \belowdisplayskip\abovedisplayskip%
  \abovedisplayshortskip\ZeroUnit%
  \belowdisplayshortskip 6\Unit%
}%
\let\OldSmall\small%
\def\small{%
  \OldSmall%
  \abovedisplayskip 8.5\Unit%
  \belowdisplayskip\abovedisplayskip%
  \abovedisplayshortskip\ZeroUnit%
  \belowdisplayshortskip 4\Unit%
}%
%
% Texts:
\def\MathNotes{Math Notes}%
%
% Title pages and covers:
\newif\ifcover%
\coverfalse%
%   - cover/title page
\def\CoverOrTitle{%
  \begingroup%
    \fullpage%
    \textbox\Zero{\lxiisf\@title}%
    \dimen\Zero\hsize%
    \advance\dimen\Zero-\wd\Zero%
    \dimen\One\FontQuadWidth\lxiisf%
    \vbox to\vsize{%
      \ifcover%
        \setcolor cmyk(.5, .5, 0, 0);%
        \rlap{\vrule width\hsize height\vsize depth\ZeroUnit}%
        \kern -.75\vsize%
      \else%
        \hrule width\ZeroUnit height\ZeroUnit depth\ZeroUnit%
        \kern .25\vsize%
      \fi%
      \moveright .5\dimen\Zero\vbox{%
        \ifcover\setcolor gray(1);\fi%
        \hrule width\wd\Zero height 1.6\Unit depth\ZeroUnit%
        \kern 5.6\Unit%
        \hrule width\wd\Zero height 4.8\Unit depth\ZeroUnit%
        \kern .125\dimen\One%
        \copy\Zero%
        \kern 2\dp\Zero%
        \kern .125\dimen\One%
        \hrule width\wd\Zero height 4.8\Unit depth\ZeroUnit%
        \kern 5.6\Unit%
        \hrule width\wd\Zero height 1.6\Unit depth\ZeroUnit%
        \kern .125\dimen\One%
        \hbox to\wd\Zero{\hss\xxvit from \@subject}%
        \ifcover\resetcolor\fi%
      }%
      \vfill%
      \ifcover\else%
        \moveright .25\dimen\Zero\vbox{%
          \dimen\Zero .5\dimen\Zero%
          \advance\dimen\Zero \wd\Zero%
          \hrule width\dimen\Zero height 1.6\Unit depth\ZeroUnit%
          \kern .125\dimen\One%
          \hbox to\dimen\Zero{%
            \input macros/t1%
            \hss\xxiiiit\@author\hss%
          }%
          \kern .05\vsize%
        }%
      \fi%
      \ifcover\resetcolor\fi%
    }%
    \eject%
  \endgroup%
}%
%   - front cover
\def\frontcover{%
  \begingroup%
    \covertrue%
    \CoverOrTitle%
  \endgroup%
}%
%   - title page
\def\titlepage{%
  \begingroup%
    \coverfalse%
    \CoverOrTitle%
  \endgroup%
}%
%   - info page
\def\infopage{%
  \vbox to\vsize{%
    \parindent\ZeroUnit%
    \parskip 12\Unit plus \Unit\relax%

    \vss%

    Title: {\it\MathNotes{} from \@subject}.

    Cover design by \@author.

    The cover and the content of this work was typeset in \TeX.

    \TeX{} is a trademark of the American Mathematical Society.

    This work is licensed under CC BY-SA 4.0.

    Copyright \copyright{} \number\year{} by \@author. You are free to copy,
    distribute, transmit and/or adapt the work under the following conditions:
    (1) You must give appropriate credit, provide a link to the license, and
    indicate if changes were made. You may do so in any reasonable manner, but
    not in any way that suggests the licensor endorses you or your use. (2) If
    you remix, transform, or build upon the material, you must distribute your
    contributions under the same or compatible license as the original.
  }%
  \eject%
}%
%   - about page
\def\aboutpage{%
  \chapter *About \MathNotes.

  The idea behind \emph{\MathNotes} originates from the need to have a compact
  collection of definitions, examples, and theorems with commented proofs from
  various topics of mathematics.

  The purpose of this work is not to be \textquotedblleft yet another
  textbook\textquotedblright{} but rather, as its name suggests, a set of notes
  taken from various sources, related to a specific topic, converted to a
  digital form. This is also reflected in organization and writing style: a
  document dedicated to a specific area groups together contributions
  (\emph{notes}) that are further grouped into chapters and sections reflecting
  specific sub-areas. Each contribution (note) tries to stay focused on the
  topic and be clear as much as possible: definitions are backed by examples,
  theorems are build up from the ground and their proofs are tried to be
  written in clean, understandable form with difficult parts explained.
  Therefore, the intended audience is a general public with interest in
  mathematics rather than matter experts which may found some parts of
  contributions trivial and explained in too much detail.
}%
%
% Setup:
\restorecat\@%
\normal%
%
% End of notes.tex:
\endinput%
